\documentclass[11pt]{article}

\usepackage[a4paper,margin=1in]{geometry}
\usepackage{amsmath,amssymb,amsthm,mathtools}
\usepackage{microtype}
\usepackage[hidelinks]{hyperref}

\newtheorem{theorem}{Theorem}[section]
\newtheorem{lemma}[theorem]{Lemma}
\newtheorem{proposition}[theorem]{Proposition}
\newtheorem{corollary}[theorem]{Corollary}
\theoremstyle{remark}
\newtheorem{remark}[theorem]{Remark}

\newcommand{\E}{\mathbb{E}}
\newcommand{\Prob}{\mathbb{P}}
\newcommand{\one}{\mathbf{1}}
\newcommand{\cost}{\operatorname{cost}}
\newcommand{\optpost}{\operatorname{OPT}_{\mathrm{post}}}
\newcommand{\optadapt}{\operatorname{OPT}_{\mathrm{adapt}}}
\newcommand{\optprior}{\operatorname{OPT}_{\mathrm{prior}}}

\title{A Probabilistic Construction for a Factor-Two\\
Lower Bound in Symmetric Stochastic TSP}
\author{}
\date{}

\begin{document}
\maketitle

\begin{abstract}
We use the probabilistic method to construct symmetric stochastic traveling
salesperson instances with a priori-to-adaptive cost ratio tending to~$2$.
The outer randomness selects a deterministic metric, whereas the inner
randomness is the independent activation process of the resulting stochastic
TSP instance.  A weighted concentration argument shows that, with high
probability over the sampled metric, every a priori master tour has expected
cost close to twice the minimum possible baseline.  A metric-dependent
adaptive policy has expected cost close to that baseline.  Consequently, some
deterministic metric in the support of the random construction has adaptivity
ratio arbitrarily close to~$2$.  The same family gives an obliviousness-gap
lower bound of~$2$, but it does not produce a nontrivial clairvoyance gap.
\end{abstract}

\section{The stochastic TSP benchmarks}

Let $V$ be a finite set of clients, let $r\notin V$ be a depot, and let
$d$ be a symmetric metric on $\{r\}\cup V$.  Each client $v\in V$ is active
independently with probability $p_v$.  The random active set is denoted by
$A\subseteq V$.

An \emph{a posteriori solution} knows $A$ before choosing a tour.  Its optimum
is
\[
    \optpost
    :=
    \E_A\bigl[\operatorname{TSP}_d(\{r\}\cup A)\bigr],
\]
where the cost for $A=\varnothing$ is defined to be~$0$.

An \emph{a priori tour} is a fixed tour
\[
    T=(r,v_1,\ldots,v_{|V|},r)
\]
on the full vertex set.  For a realization $A$, let $T|_A$ be the tour obtained
by retaining only $r$ and the active clients, in their order in~$T$.  Its
expected cost is
\[
    \cost_d(T):=\E_A\bigl[\cost_d(T|_A)\bigr],
\]
and
\[
    \optprior(d):=\min_T \cost_d(T).
\]

An \emph{adaptive policy} calls uncalled clients one at a time.  If the called
client is active, the salesperson must travel to it immediately; if it is
inactive, the salesperson remains at the current location.  The next call may
depend on the current location and on all outcomes observed so far.  After all
clients have been called, the salesperson returns to~$r$.  For a policy~$\pi$,
write
\[
    \cost_d(\pi):=\E_A\bigl[\cost_d(\pi,A)\bigr],
    \qquad
    \optadapt(d):=\min_\pi \cost_d(\pi).
\]
The metric and all activation probabilities are known to the policy; only the
realization $A$ is initially unknown.

For every fixed instance,
\[
    \optpost(d)\leq \optadapt(d)\leq \optprior(d).
\]
The \emph{adaptivity ratio} is $\optprior/\optadapt$, and the
\emph{obliviousness ratio} is $\optprior/\optpost$.

\section{Outer randomness and the probabilistic method}

The random metric used below is only a proof device.  There are two independent
sources of randomness:
\[
    \underbrace{D}_{\text{randomly sampled metric}}
    \qquad\text{and}\qquad
    \underbrace{A}_{\text{active set in the fixed instance}}.
\]
After a realization $D=d$ is selected, the stochastic TSP instance is the
ordinary deterministic metric instance $(\{r\}\cup V,d,\{p_v\})$.  Both the
optimal a priori tour and the optimal adaptive policy may depend on~$d$.
Thus the sampled metric is not hidden information available to one benchmark
but not another.

The following elementary statement records the relevant probabilistic-method
principle.

\begin{proposition}[Existence from random certificates]\label{prop:method}
Let $I_D$ be a random deterministic instance, and let $X(I)$ and $Y(I)$ be two
nonnegative optimum values.  Suppose that, for every realization~$I$,
\[
    X(I)\geq L(I)
    \qquad\text{and}\qquad
    Y(I)\leq U(I).
\]
If
\[
    \Prob_D\bigl[L(I_D)\geq \rho U(I_D)\bigr]>0,
\]
then there is a deterministic realization~$I$ satisfying
$X(I)/Y(I)\geq\rho$.  The same conclusion follows from
\[
    \E_D\bigl[L(I_D)-\rho U(I_D)\bigr]\geq 0.
\]
\end{proposition}

\begin{proof}
In the first case, choose any realization in the positive-probability event.
Then
\[
    X(I)\geq L(I)\geq \rho U(I)\geq \rho Y(I).
\]
In the second case, some realization satisfies
$L(I)-\rho U(I)\geq0$, and the same chain of inequalities applies.
\end{proof}

A central quantifier issue is that
\[
    \E_D\!\left[\min_T \cost_D(T)\right]
    \leq
    \min_T \E_D\!\left[\cost_D(T)\right].
\]
Therefore, proving that every \emph{fixed} tour is expensive on average over
$D$ does not suffice: the optimizing tour is allowed to depend on the sampled
metric.  Our proof establishes a simultaneous high-probability lower bound for
all tours.

\section{The random metric construction}

All logarithms are natural.  Fix a sufficiently large integer~$n$, let
$V=[n]$, and set
\[
    p=p_n:=\frac{1}{(\log n)^4},
    \qquad
    \theta=\theta_n:=\frac{1}{\log n},
    \qquad
    q:=1-p.
\]
Every client is active independently with probability~$p$.

Independently for every unordered pair
$e=\{x,y\}\subseteq\{r\}\cup V$, sample $X_e\sim\operatorname{Bernoulli}(\theta)$
and define
\[
    D(x,y):=2-X_{\{x,y\}}
    \quad (x\neq y),
    \qquad
    D(x,x):=0.
\]
Thus every pair has distance~$1$ with probability~$\theta$ and distance~$2$
with probability~$1-\theta$, independently of all other pairs.  We call a
pair of distance~$1$ a \emph{cheap edge}.

Every realization of~$D$ is a symmetric metric.  Indeed, for distinct
$x,y,z$,
\[
    D(x,z)\leq 2\leq D(x,y)+D(y,z),
\]
and all remaining triangle inequalities are immediate.

Let $K:=|A|\sim\operatorname{Binomial}(n,p)$ and define
\[
    W=W_n
    :=\E_A\bigl[(K+1)\one_{\{K\geq1\}}\bigr]
    =np+1-q^n.
\]
For a nonempty realization with $K=k$, every tour through $r$ and the active
clients has exactly $k+1$ traveled edges.  Since every distance is at least~$1$,
$W$ is a universal lower bound on the expected cost of every feasible
benchmark.  Moreover,
\[
    W\geq np=\frac{n}{(\log n)^4},
    \qquad
    \theta W\geq\frac{n}{(\log n)^5}\longrightarrow\infty.
\]

\section{A simultaneous lower bound for all a priori tours}

Fix an a priori tour
\[
    T=(r,v_1,v_2,\ldots,v_n,r)
\]
and a deterministic realization~$d$ of the metric.  Write
\[
    C_d(T):=\E_A\bigl[\cost_d(T|_A)\bigr].
\]

\begin{lemma}[Exact expected-cost expansion]\label{lem:expansion}
For $i\in[n]$, define
\[
    \alpha_i:=p\bigl(q^{i-1}+q^{n-i}\bigr),
\]
and for $1\leq i<j\leq n$, define
\[
    \beta_{ij}:=p^2q^{j-i-1}.
\]
Then
\begin{equation}\label{eq:cost-expansion}
    C_d(T)
    =
    \sum_{i=1}^n \alpha_i d(r,v_i)
    +
    \sum_{1\leq i<j\leq n}\beta_{ij}d(v_i,v_j).
\end{equation}
Furthermore, the sum of all coefficients in \eqref{eq:cost-expansion} is~$W$.
\end{lemma}

\begin{proof}
The edge from $r$ to $v_i$ is used when $v_i$ is active and
$v_1,\ldots,v_{i-1}$ are inactive, an event of probability $pq^{i-1}$.
The edge from $v_i$ to $r$ is used when $v_i$ is active and
$v_{i+1},\ldots,v_n$ are inactive, an event of probability $pq^{n-i}$.
For $i<j$, the clients $v_i$ and $v_j$ are consecutive active clients in the
shortcut whenever they are both active and every client strictly between them
is inactive.  This event has probability $p^2q^{j-i-1}$.  This proves
\eqref{eq:cost-expansion}.

If every distance in \eqref{eq:cost-expansion} is replaced by~$1$, the right
side becomes the expected number of traveled edges in the shortcut tour.  That
number is $(K+1)\one_{\{K\geq1\}}$, whose expectation is~$W$.  Hence the
coefficients sum to~$W$.
\end{proof}

We use the following weighted form of Hoeffding's inequality: if
$Y_1,\ldots,Y_m$ are independent random variables taking values in $[0,1]$ and
$a_1,\ldots,a_m\geq0$, then, for every $t>0$,
\begin{equation}\label{eq:weighted-hoeffding}
    \Prob\!\left[
        \sum_{\ell=1}^m a_\ell
        \bigl(Y_\ell-\E Y_\ell\bigr)\geq t
    \right]
    \leq
    \exp\!\left(-\frac{2t^2}{\sum_{\ell=1}^m a_\ell^2}\right).
\end{equation}

\begin{lemma}[Coefficient-square bound]\label{lem:squares}
For every a priori order and all sufficiently large~$n$, the coefficients in
\eqref{eq:cost-expansion} satisfy
\[
    \sum_{i=1}^n\alpha_i^2
    +
    \sum_{1\leq i<j\leq n}\beta_{ij}^2
    \leq 2np^3.
\]
\end{lemma}

\begin{proof}
Using $(a+b)^2\leq2a^2+2b^2$,
\begin{align*}
    \sum_{i=1}^n\alpha_i^2
    &\leq
    2p^2\sum_{i=1}^n
        \bigl(q^{2(i-1)}+q^{2(n-i)}\bigr)\\
    &\leq
    \frac{4p^2}{1-q^2}
    =\frac{4p}{2-p}
    \leq 4p.
\end{align*}
Also,
\begin{align*}
    \sum_{1\leq i<j\leq n}\beta_{ij}^2
    &=p^4\sum_{g=0}^{n-2}(n-g-1)q^{2g}\\
    &\leq
    \frac{np^4}{1-q^2}
    =\frac{np^3}{2-p}
    \leq np^3.
\end{align*}
Since $np^2=n/(\log n)^8\to\infty$, we have $4p\leq np^3$ for all
sufficiently large~$n$.  Adding the two bounds proves the claim.
\end{proof}

\begin{lemma}[All master tours are expensive]\label{lem:all-tours}
With probability $1-o(1)$ over the random metric~$D$,
\[
    \optprior(D)
    \geq
    \left(2-\frac{2}{\log n}\right)W.
\]
\end{lemma}

\begin{proof}
Fix a tour~$T$.  Let $c_e$ denote the coefficient of the unordered pair~$e$ in
\eqref{eq:cost-expansion}; thus $c_{\{r,v_i\}}=\alpha_i$ and
$c_{\{v_i,v_j\}}=\beta_{ij}$.  Since $D(e)=2-X_e$,
\[
    C_D(T)=2W-\sum_e c_eX_e.
\]
By Lemma~\ref{lem:expansion}, $\sum_e c_e=W$, and therefore
\[
    \E_D[C_D(T)]=(2-\theta)W.
\]
Set
\[
    \varepsilon:=\frac{1}{\log n}.
\]
The event
$C_D(T)\leq(2-\theta-\varepsilon)W$ is the event
\[
    \sum_e c_e(X_e-\theta)\geq\varepsilon W.
\]
Applying \eqref{eq:weighted-hoeffding}, Lemma~\ref{lem:squares}, and
$W\geq np$ gives
\begin{align*}
    \Prob_D\!\left[
        C_D(T)\leq(2-\theta-\varepsilon)W
    \right]
    &\leq
    \exp\!\left(
        -\frac{2\varepsilon^2W^2}{\sum_e c_e^2}
    \right)\\
    &\leq
    \exp\!\left(-\frac{\varepsilon^2W^2}{np^3}\right)\\
    &\leq
    \exp\!\left(-\frac{\varepsilon^2n}{p}\right)\\
    &=\exp\!\bigl(-n(\log n)^2\bigr).
\end{align*}
There are at most $n!$ a priori tours with the depot fixed.  Hence, by a union
bound,
\begin{align*}
    &\Prob_D\!\left[
        \text{some tour }T\text{ satisfies }
        C_D(T)<(2-\theta-\varepsilon)W
    \right]\\
    &\hspace{3cm}\leq
    n!\exp\!\bigl(-n(\log n)^2\bigr)\\
    &\hspace{3cm}\leq
    \exp\!\bigl(n\log n-n(\log n)^2\bigr)
    =o(1).
\end{align*}
Thus, with probability $1-o(1)$, the stated lower bound holds simultaneously
for every a priori tour.  Since $\theta=\varepsilon=1/\log n$, the claimed
factor is $2-2/\log n$.
\end{proof}

\section{A cheap adaptive policy}

For each deterministic metric realization~$d$, define a deterministic adaptive
policy~$\pi_d$ as follows.  Suppose the salesperson is at $u$ and the set of
uncalled clients is~$S$.
\begin{itemize}
    \item If $S=\varnothing$, return to~$r$.
    \item If there exists $v\in S$ with $d(u,v)=1$, call the least-labeled such
    vertex.
    \item Otherwise, call the least-labeled vertex of~$S$.
\end{itemize}
Thus, while the salesperson remains at a fixed location~$u$, all uncalled cheap
neighbors of~$u$ are called before any distance-$2$ vertex is called.

The policy is allowed to depend on~$d$, because the metric is part of the fixed
instance and is known before the activation outcomes are observed.

\begin{lemma}[Cost of the adaptive policy, averaged over metrics]
\label{lem:adaptive-average}
Let
\[
    U(d):=\E_A\bigl[\cost_d(\pi_d,A)\bigr].
\]
Then
\[
    W\leq U(d)\quad\text{for every }d,
\]
and
\[
    \E_D[U(D)]
    \leq
    np+\left(\frac{1}{\theta}+1\right)(1-q^n)
    \leq W+\frac{1}{\theta}.
\]
\end{lemma}

\begin{proof}
First fix $k\geq1$ and condition on $K=k$.  Consider the execution jointly over
the random metric and the active set.  Immediately after the salesperson
arrives at a new active client~$u$, let $S$ be the set of uncalled clients and
suppose that exactly $s$ clients in~$S$ are active.

We may expose the random metric lazily.  Before this arrival, none of the edges
$\{u,v\}$ with $v\in S$ was needed by the policy.  Indeed, prior choices use
only edges incident to the then-current vertex, and $u$ had not previously been
the current vertex.  The only previously exposed edge incident to~$u$ may be
the edge from the preceding current vertex to~$u$; that preceding vertex is not
in~$S$.  At the initial state $u=r$, no edge incident to an uncalled client has
yet been exposed.  Consequently, conditional on the entire preceding history,
the edges $\{u,v\}$ for $v\in S$ remain independent, each being cheap with
probability~$\theta$, and they remain independent of the activation states of
the clients in~$S$.

The next traveled edge has length~$2$ exactly when none of the $s$ remaining
active clients is a cheap neighbor of~$u$.  The probability of this event is
$(1-\theta)^s$.  Hence the conditional expected length of the next traveled
edge is
\[
    1+(1-\theta)^s.
\]
As the policy visits the $k$ active clients, the number~$s$ of active clients
still unvisited takes the values $k,k-1,\ldots,1$.  The final return to~$r$ costs
at most~$2$.  Therefore
\begin{align*}
    \E_{D,A}\bigl[\cost_D(\pi_D,A)\mid K=k\bigr]
    &\leq
    \sum_{s=1}^k\bigl(1+(1-\theta)^s\bigr)+2\\
    &\leq
    k+\frac{1-\theta}{\theta}+2\\
    &=k+\frac{1}{\theta}+1.
\end{align*}
For $K=0$, the cost is~$0$.  Averaging over~$K$ yields
\begin{align*}
    \E_D[U(D)]
    &=\E_{D,A}\bigl[\cost_D(\pi_D,A)\bigr]\\
    &\leq
    \E[K]
    +\left(\frac{1}{\theta}+1\right)\Prob[K\geq1]\\
    &=np+\left(\frac{1}{\theta}+1\right)(1-q^n)\\
    &=W+\frac{1-q^n}{\theta}
    \leq W+\frac{1}{\theta}.
\end{align*}

For the pointwise lower bound, if $K=k\geq1$, the policy travels along exactly
$k+1$ edges: one edge to each active client and one final edge back to~$r$.
Every such edge has length at least~$1$.  Thus its realized cost is at least
$(k+1)\one_{\{k\geq1\}}$, and taking expectation over~$A$ gives $U(d)\geq W$.
\end{proof}

\begin{lemma}[A good adaptive policy for most sampled metrics]
\label{lem:adaptive-whp}
Define
\[
    \delta_n:=\frac{1}{\sqrt{\theta W}}.
\]
Then $\delta_n=o(1)$ and
\[
    \Prob_D\bigl[\optadapt(D)\leq(1+\delta_n)W\bigr]
    \geq 1-\delta_n.
\]
\end{lemma}

\begin{proof}
Since $\theta W\to\infty$, we have $\delta_n=o(1)$.  By
Lemma~\ref{lem:adaptive-average}, the nonnegative random variable $U(D)-W$
satisfies
\[
    \E_D[U(D)-W]\leq\frac{1}{\theta}.
\]
Markov's inequality therefore gives
\begin{align*}
    \Prob_D\bigl[U(D)>(1+\delta_n)W\bigr]
    &=
    \Prob_D\bigl[U(D)-W>\delta_nW\bigr]\\
    &\leq
    \frac{1/\theta}{\delta_nW}
    =\frac{1}{\sqrt{\theta W}}
    =\delta_n.
\end{align*}
Since $\optadapt(d)\leq U(d)$ for every~$d$, the claim follows.
\end{proof}

\section{Extracting a deterministic metric}

\begin{theorem}[Factor-two adaptivity lower bound]\label{thm:main}
There exists a sequence of deterministic symmetric stochastic TSP instances
for which
\[
    \frac{\optprior}{\optadapt}\longrightarrow 2.
\]
Consequently, the adaptivity gap for symmetric stochastic TSP is at least~$2$.
\end{theorem}

\begin{proof}
By Lemma~\ref{lem:all-tours}, with probability $1-o(1)$,
\[
    \optprior(D)\geq
    \left(2-\frac{2}{\log n}\right)W.
\]
By Lemma~\ref{lem:adaptive-whp}, with probability $1-o(1)$,
\[
    \optadapt(D)\leq(1+\delta_n)W.
\]
The two events therefore occur simultaneously with positive probability for all
sufficiently large~$n$.  Choose one deterministic realization~$d_n$ in their
intersection.  Then
\[
    \frac{\optprior(d_n)}{\optadapt(d_n)}
    \geq
    \frac{2-2/\log n}{1+\delta_n}.
\]
Since $\delta_n=o(1)$, the right side tends to~$2$.
\end{proof}

The proof does not identify a particular successful realization~$d_n$.  It
shows more: a random realization succeeds with probability tending to~$1$.
Thus the deterministic metric is obtained nonconstructively by the
probabilistic method.

\begin{corollary}[Factor-two obliviousness lower bound]
There exists a sequence of deterministic symmetric stochastic TSP instances
for which
\[
    \frac{\optprior}{\optpost}\geq 2-o(1).
\]
In particular, the obliviousness gap is at least~$2$.
\end{corollary}

\begin{proof}
For every instance, $\optpost\leq\optadapt$.  Hence
\[
    \frac{\optprior}{\optpost}
    \geq
    \frac{\optprior}{\optadapt},
\]
and Theorem~\ref{thm:main} applies.
\end{proof}

\section{Why this construction has no clairvoyance gap}

For every realization~$d$ of the metric,
\[
    \optpost(d)\geq W,
\]
because a nonempty active set of size~$k$ requires a tour with $k+1$ edges and
every edge has length at least~$1$.  For the deterministic realizations selected
in the proof of Theorem~\ref{thm:main},
\[
    W\leq\optpost(d_n)\leq\optadapt(d_n)\leq(1+\delta_n)W.
\]
Consequently,
\[
    1\leq
    \frac{\optadapt(d_n)}{\optpost(d_n)}
    \leq 1+\delta_n
    =1+o(1).
\]
Thus random $\{1,2\}$-metrics separate the a priori and adaptive benchmarks,
but they do not separate the adaptive and a posteriori benchmarks.

\section{A Bellman certificate for future clairvoyance constructions}

This section records a general device for a different probabilistic
construction; it does not itself prove a new clairvoyance lower bound.

For a fixed metric~$d$ and independent activation probabilities $p_v$, let
$F_d(u,S)$ be the optimal expected remaining cost of an adaptive policy when the
salesperson is at~$u$ and the set of uncalled clients is~$S$.  Then
\[
    F_d(u,\varnothing)=d(u,r),
\]
and, for $S\neq\varnothing$,
\begin{equation}\label{eq:bellman}
    F_d(u,S)
    =
    \min_{v\in S}
    \left\{
        (1-p_v)F_d(u,S\setminus\{v\})
        +p_v\bigl(d(u,v)+F_d(v,S\setminus\{v\})\bigr)
    \right\}.
\end{equation}
In particular, $\optadapt(d)=F_d(r,V)$.

A function $\Phi_d(u,S)$ is a lower-bound certificate if
\[
    \Phi_d(u,\varnothing)\leq d(u,r)
\]
and, for every nonempty $S$ and every $v\in S$,
\begin{equation}\label{eq:subsolution}
    \Phi_d(u,S)
    \leq
    (1-p_v)\Phi_d(u,S\setminus\{v\})
    +p_v\bigl(d(u,v)+\Phi_d(v,S\setminus\{v\})\bigr).
\end{equation}
Induction on $|S|$, using \eqref{eq:bellman}, gives
\[
    \Phi_d(u,S)\leq F_d(u,S)
\]
for every state, and hence
\[
    \Phi_d(r,V)\leq\optadapt(d).
\]
This avoids a direct union bound over adaptive policies.  In a probabilistic
clairvoyance-gap construction, one may try to show that a random metric admits
both
\begin{enumerate}
    \item a low-cost a posteriori routing for most activation realizations, and
    \item a Bellman subsolution $\Phi_d$ that forces every adaptive policy to
    incur a larger expected cost.
\end{enumerate}
The random structure would need to couple many routing decisions whose correct
resolution depends on activation information not yet revealed to the adaptive
policy.

\end{document}
